\begin{table*}[!htbp]
\centering
\scriptsize
\caption{Non-Markov Brownian bridge stress}
\label{tab:b2}
\vspace{0.5em}
\par\textbf{Panel A. Coverage and width}

\begin{adjustbox}{max width=\linewidth}
\begin{tabular}{@{}lllllllll@{}}
\toprule
Environment & Tier & open-only cov. (SE) & open width & cross-event cov. (SE) & cross width & width ratio & $\tau$-stratified cov. (SE) & $\tau$-stratified width ratio \\
\midrule
$\nu_{\rm mult}=0$ placebo & oracle & 0.897 (0.008) & 0.1618 & 0.897 (0.008) & 0.1637 & 1.01$\times$ & 0.901 (0.008) & 1.01$\times$ \\
$\nu_{\rm mult}=0$ placebo & estimated & 0.882 (0.008) & 0.1576 & 0.900 (0.008) & 0.1662 & 1.05$\times$ & 0.900 (0.008) & 1.05$\times$ \\
$\nu_{\rm mult}=0.6$ bridge & oracle & 0.818 (0.010) & 0.1618 & 0.900 (0.008) & 0.2047 & 1.27$\times$ & 0.902 (0.008) & 1.27$\times$ \\
$\nu_{\rm mult}=0.6$ bridge & estimated & 0.806 (0.010) & 0.1571 & 0.900 (0.008) & 0.2072 & 1.32$\times$ & 0.901 (0.008) & 1.31$\times$ \\
\bottomrule
\end{tabular}
\end{adjustbox}
\vspace{0.5em}
\par\noindent\begin{minipage}{\linewidth}\raggedright{\scriptsize Notes. 90\% prediction intervals, 5-fold out-of-fold (fold = seed mod 5), 128 seeds $\times$ 4 windows = 512 paths, 1,536 double-endpoint events ($\tau=2/65/81$, 512 each). Nominal 0.90 ($\alpha=0.1$). Coverage SEs are binomial over the 1,536 event predictions and ignore within-path and across-fold dependence, so they are lower bounds. \texttt{open-only} uses ordinary OPEN$\to$OPEN one-step transitions and the Markov $\tau$-step variance $v_{0e}$; \texttt{cross-event} uses $\operatorname{Var}(r_e)=v_{0e}+\theta\sigma_e^2\,\mathrm{gap}_e$ with $\theta$ fitted by Gaussian MLE; \texttt{$\tau$-stratified} uses a separate conformal quantile per freeze length and no functional form in $\tau$.}\end{minipage}\par

\par\textbf{Panel B. Bridge variance component}

\begin{adjustbox}{max width=\linewidth}
\begin{tabular}{@{}llllll@{}}
\toprule
Environment & Tier & $\hat\theta$ & 95\% CI & $\hat\nu/\sigma$ & 95\% CI \\
\midrule
$\nu_{\rm mult}=0$ placebo & oracle & 0.0282 & [0.000, 0.097] & 0.168 & [0.000, 0.311] \\
$\nu_{\rm mult}=0$ placebo & estimated & 0.0703 & [0.008, 0.138] & 0.265 & [0.092, 0.372] \\
$\nu_{\rm mult}=0.6$ bridge & oracle & 0.4912 & [0.397, 0.595] & 0.701 & [0.630, 0.771] \\
$\nu_{\rm mult}=0.6$ bridge & estimated & 0.5360 & [0.443, 0.639] & 0.732 & [0.665, 0.799] \\
\bottomrule
\end{tabular}
\end{adjustbox}
\vspace{0.5em}
\par\noindent\begin{minipage}{\linewidth}\raggedright{\scriptsize Notes. $\theta=\nu_{\rm mult}^2\,\mathbb E[e^{2\zeta}]$; the design value is $\theta_{\rm true}=0.457$ from the realized event-level draws. The $\nu=0$ placebo has $\theta=0$ by construction. The oracle-tier placebo interval contains zero; the estimated-tier placebo has a small positive bias from fitting $T$ and $\sigma$ on open transitions, as in the $\gamma_{\mathrm{st}}=0$ placebo of Appendix Table B10.}\end{minipage}\par

\par\textbf{Panel C. Coverage by freeze length}

\begin{adjustbox}{max width=\linewidth}
\begin{tabular}{@{}llllllllll@{}}
\toprule
Environment & Tier & $\tau$ & $n$ & open cov. & open width & cross cov. & cross width & $\tau$-strat cov. & $\tau$-strat width \\
\midrule
$\nu_{\rm mult}=0$ placebo & oracle & 2 & 512 & 0.891 & 0.0488 & 0.893 & 0.0492 & 0.900 & 0.0497 \\
$\nu_{\rm mult}=0$ placebo & oracle & 65 & 512 & 0.891 & 0.2089 & 0.891 & 0.2113 & 0.902 & 0.2188 \\
$\nu_{\rm mult}=0$ placebo & oracle & 81 & 512 & 0.910 & 0.2277 & 0.908 & 0.2305 & 0.900 & 0.2239 \\
$\nu_{\rm mult}=0$ placebo & estimated & 2 & 512 & 0.887 & 0.0479 & 0.902 & 0.0502 & 0.898 & 0.0499 \\
$\nu_{\rm mult}=0$ placebo & estimated & 65 & 512 & 0.865 & 0.2034 & 0.893 & 0.2145 & 0.902 & 0.2213 \\
$\nu_{\rm mult}=0$ placebo & estimated & 81 & 512 & 0.895 & 0.2213 & 0.906 & 0.2339 & 0.900 & 0.2273 \\
$\nu_{\rm mult}=0.6$ bridge & oracle & 2 & 512 & 0.820 & 0.0488 & 0.895 & 0.0597 & 0.904 & 0.0612 \\
$\nu_{\rm mult}=0.6$ bridge & oracle & 65 & 512 & 0.797 & 0.2089 & 0.889 & 0.2641 & 0.900 & 0.2800 \\
$\nu_{\rm mult}=0.6$ bridge & oracle & 81 & 512 & 0.838 & 0.2277 & 0.916 & 0.2903 & 0.900 & 0.2744 \\
$\nu_{\rm mult}=0.6$ bridge & estimated & 2 & 512 & 0.812 & 0.0479 & 0.902 & 0.0607 & 0.900 & 0.0610 \\
$\nu_{\rm mult}=0.6$ bridge & estimated & 65 & 512 & 0.791 & 0.2029 & 0.879 & 0.2672 & 0.902 & 0.2879 \\
$\nu_{\rm mult}=0.6$ bridge & estimated & 81 & 512 & 0.814 & 0.2206 & 0.920 & 0.2935 & 0.900 & 0.2708 \\
\bottomrule
\end{tabular}
\end{adjustbox}
\vspace{0.5em}
\par\noindent\begin{minipage}{\linewidth}\raggedright{\scriptsize Notes. The open-only deficit is present at every $\tau$; it is not a single long-freeze cell. The $\tau$-stratified calibrator recovers within-$\tau$ marginal coverage without assuming the bridge variance shape.}\end{minipage}\par

\par\textbf{Panel D. Same endpoints, different bridge volatility}

\begin{adjustbox}{max width=\linewidth}
\begin{tabular}{@{}ll@{}}
\toprule
Quantity & Value \\
\midrule
Same-endpoint events & 64 \\
Max abs pool-price difference & 0.0 \\
Observed paths bit-identical & True \\
Mean latent QV at $\nu_e$ & 1.6023e-03 \\
Mean latent QV at $2\nu_e$ & 6.1677e-03 \\
Empirical QV ratio & 3.85 \\
Theoretical QV ratio & 4.00 \\
\bottomrule
\end{tabular}
\end{adjustbox}
\vspace{0.5em}
\par\noindent\begin{minipage}{\linewidth}\raggedright{\scriptsize Notes. Each selected event is re-simulated twice with the same realized endpoints, the same standard Brownian bridge shape, and volatilities $\nu_e$ and $2\nu_e$. Freeze interiors are masked by the observation map, so the observed pool-price paths are bit-identical; only the latent quadratic variation changes. The empirical QV ratio is close to the theoretical $2^2=4$. This is the empirical counterpart of the Theorem 4 construction, not a coverage experiment: with endpoints pinned, no observable-data estimator can distinguish the two bridge laws.}\end{minipage}\par

\par\textbf{Panel E. Coverage against the realized event-level volatility}

\begin{adjustbox}{max width=\linewidth}
\begin{tabular}{@{}llllllll@{}}
\toprule
$\nu_e$ quartile & $n$ & $\nu$ min & $\nu$ max & open cov. & cross cov. & open width & cross width \\
\midrule
0 & 384 & 0.0010 & 0.0034 & 0.859 & 0.919 & 0.1210 & 0.1534 \\
1 & 384 & 0.0034 & 0.0048 & 0.844 & 0.917 & 0.1451 & 0.1839 \\
2 & 384 & 0.0048 & 0.0064 & 0.797 & 0.904 & 0.1654 & 0.2093 \\
3 & 384 & 0.0064 & 0.0293 & 0.773 & 0.859 & 0.2157 & 0.2724 \\
\bottomrule
\end{tabular}
\end{adjustbox}
\vspace{0.5em}
\par\noindent\begin{minipage}{\linewidth}\raggedright{\scriptsize Notes. The groups use the simulation-only $\nu_e$, which is not observed by either method; they are a dose-response diagnostic, not a deployable stratification. Open-only coverage falls monotonically across all four quartiles; cross-event calibration recovers marginal coverage but not coverage conditional on the latent $\nu_e$ (the highest quartile remains below nominal). The latter is expected: $\nu_e$ is not in the observed information set.}\end{minipage}\par

\par\noindent\begin{minipage}{\linewidth}\raggedright{\scriptsize Bridge law Brownian bridge pinned at the latent pre/post anchors. The post endpoint is drawn from $N(m_+, v_{\rm ar}+\nu_e^2\,\mathrm{gap})$, with $\nu_e=\sigma\nu_{\rm mult}e^{\zeta_e}$, $\zeta_e\sim N(0,0.35^2)$ clipped to $[-1.0,1.0]$, and $\nu_{\rm mult}=0.60$. The interior is unobserved, so the bridge volatility's observable footprint is the endpoint law; the same-endpoint panel isolates its path-functional consequence.}\end{minipage}\par

\end{table*}
